\indent ಪದಗಳನ್ನು ಸಾಮಾನ್ಯವಾಗಿ ಬಳಸುತ್ತೇವೆ. ಇದೇ ರೀತಿಯಾಗಿ ಎರಡು ವಾಕ್ಯಗಳನ್ನು ಜೋಡಿಸಲು “ಮತ್ತು, ಆದರೆ, ಅಥವಾ, ಆದ್ದರಿಂದ, ಹಾಗಾದರೆ, ಇನ್ನೂ, ಆದುದರಿಂದ” ಇತ್ಯಾದಿ ಪದಗಳನ್ನು ಬಳಸುತ್ತೇವೆ. ಈ ಬಗೆಯ ಪದಗಳಿಗೆ ಸಂಬಂಧಾರ್ಥಕಾವ್ಯಯಗಳೆಂದು ಹೆಸರು.\par

\vspace{3pt plus 2pt}\noindent\hspace*{1.2em}\kbvlabel{ಉದಾ}\par\vspace{1pt}

\vspace{2.5pt plus 3pt minus 0.5pt}\noindent\hspace*{1.2em}\kbvnum{(೧)}\hspace{0.55em}ಹಸು ಮತ್ತು ಕರು\par

\vspace{2.5pt plus 3pt minus 0.5pt}\noindent\hspace*{1.2em}\kbvnum{(೨)}\hspace{0.55em}ಹಸುವೂ-ಕರುವೂ\par

\vspace{2.5pt plus 3pt minus 0.5pt}\noindent\hspace*{1.2em}\kbvnum{(೩)}\hspace{0.55em}ಅಕ್ಕನೂ ತಂಗಿಯೂ\par

\vspace{2.5pt plus 3pt minus 0.5pt}\noindent\hspace*{1.2em}\kbvnum{(೪)}\hspace{0.55em}ಬೆಲ್ಲ ಅಥವಾ ಸಕ್ಕರೆ\par

\vspace{2.5pt plus 3pt minus 0.5pt}\noindent\hspace*{1.2em}\kbvnum{(೫)}\hspace{0.55em}ಕುದುರೆಯೂ ಬಂದಿತು ; ಗಾಡಿಯೂ ಬಂದಿತು.\par

\vspace{3pt plus 2pt}\noindent\hspace*{1.2em}\kbvlabel{ಉದಾ}\par\vspace{1pt}

\vspace{2.5pt plus 3pt minus 0.5pt}\noindent\hspace*{1.2em}\kbvnum{(೧)}\hspace{0.55em}ಅವಳು ಆಂದ್ರಾಗಿ ತಂದಳು ಮತ್ತು ಅದನ್ನು ಬೀಸಿ ವಳ ಗಡಿಯಿಂದ ತಂದಳು ಮತ್ತು ದನ್ನು ಬಿಹಟ್ಟು ಮಾಡಿದಳು.\par

\vspace{2.5pt plus 3pt minus 0.5pt}\noindent\hspace*{1.2em}\kbvnum{(೨)}\hspace{0.55em}ರೈಲು ನಿಲ್ದಾಣಕ್ಕೆ ಬನ್ನಿ ಅಥವಾ ನೇರವಾಗಿ ಕಲ್ಯಾಣ ಮಂಟಪಕ್ಕೆ ಬನ್ನಿ.\par

\vspace{2.5pt plus 3pt minus 0.5pt}\noindent\hspace*{1.2em}\kbvnum{(೩)}\hspace{0.55em}ಮಕ್ಕಳನ್ನು ಮುದ್ದಿಸಬೇಕು ಆದರೆ ಅತಿಯಾಗಿ ಮುದ್ದಿಸುವುದು ಸರ್ವಥಾ ತಪ್ಪು.\par

\vspace{2.5pt plus 3pt minus 0.5pt}\noindent\hspace*{1.2em}\kbvnum{(೪)}\hspace{0.55em}ರೈಲು ಬಹಳ ತಡವಾಗಿ ಬಂದಿತು ಆದುದರಿಂದ ನಾನು ಪ್ರಯಾಣವನ್ನು ಮುಂದಕ್ಕೆ ಹಾಕಿದೆ.\par

\vspace{2.5pt plus 3pt minus 0.5pt}\noindent\hspace*{1.2em}\kbvnum{(೫)}\hspace{0.55em}ನೀನು ಪ್ರಯಾಣದ ಖರ್ಚು ಕೊಡಬೇಕು ಹಾಗಾದರೆ ನಾನು ನಿನ್ನೊಂದಿಗೆ ಬರುತ್ತೇನೆ.\par

\vspace{2.5pt plus 3pt minus 0.5pt}\noindent\hspace*{1.2em}\kbvnum{(೬)}\hspace{0.55em}ನೀನು ತುಂಬಾ ಸಹಾಯ ಮಾಡಿದ್ದೀಯೆ ಇನ್ನೂ ಹೆಚ್ಚಿನ ಸಹಾಯ ನಿನ್ನಿಂದ ಆಗಬೇಕಾಗಿದೆ.\par

\kbvsechead{(೬)\hspace{0.45em}ಕೃದಂತಾವ್ಯಯ :}

\kbvsubtitle{ಕೃದಂತ ಪ್ರಕರಣ ನೋಡಿ}

\kbvsechead{(೭)\hspace{0.45em}ತದ್ದಿತಾಂತಾವ್ಯಯ :}

\kbvsubtitle{ತದ್ದಿತಾಂತ ಪ್ರಕರಣ ನೋಡಿ}
